\documentclass{extreport}
\usepackage{graphicx} % Required for inserting images
\usepackage{titlesec}
\usepackage{geometry}
\usepackage{parskip}
\pagenumbering{gobble}

\geometry{
a4paper,
top=10mm,
left=10mm,
right=10mm,
bottom=10mm,
}



\title{IT Act assignments}
\author{Hrishikesh MJ}
\date{}

\begin{document}


\maketitle

% \centering \title{IT Act 2000}

\section*{Chapter I - Preliminary}

\par Chapter I titled "Preliminary" sets the foundation for the IT Act 2000. This chapter clarifies that the Act applies to the whole territory of India and also to offenses committed outside India which involves a computer or network in India. The chapter also provides definitions for key terms that will be used throughout the Act.

\section*{Chapter II - Digital Signature and Electronic Signature}

\par Chapter II of the act gives legal recognition to digital signatures and also introduces electronic signatures. This chapter also explains how electronic records can be authenticated using an electronic signature and cryptographic techniques. 

\section*{Chapter III - Electronic Governance}

\par This chapter promotes e-governance or electric governance by recognizing the validity of electronic records and signatures in official dealings. This chapter also allows the use of electronic records, retention of records. publication of rules in electronic form, and recognizes the validity of electronic contracts.

\section*{Chapter IV - Attribution, Acknowledgment and Despatch of Electronic Records}

\par This chapter sets the rules for attributing electronic records to individuals. This chapter also specifies the circumstances under which acknowledgments are required and how disputes regarding delivery and timing can be resolved. 

\section*{Chapter V - Secure Electronic Records and Secure Electronic Signature} 

\par Chapter V states that all electronic records and signatures can be considered "secure" if certain security procedures are used to authenticate the records and if the signatures are under the exclusive control of the signatory. This chapter also enables the Central Government to prescribe detailed security procedures and practices. 

\section*{Chapter VI - Regulation of Certifying Authorities}

\par This chapter deals with the licensing, regulation, and oversight of Certifying Authorities. This chapter also describes the appointment and powers of the Controller of Certifying authorities, the process of granting and revoking licenses, and the duties every Certifying Authorities have to follow to ensure secure and trustworthy operations. 

\section*{Chapter VII - Electronic Signature Certificates}

\par Chapter VII deals with the issuance and management of electronic signature certificates by Certifying Authorities, It includes the procedure for individuals or organizations to apply for such certificates. A Digital Signature Certificate is issued by a Certifying Authority only if certain conditions are met. This chapter also gives the Certifying Authority the right to suspend a Digital Signature Certificate. 

\section*{Chapter VIII - Duties of Subscribers}

\par Chapter VIII of the IT Act imposes specific duties on subscribers (Individuals or entities who are issued electronic signature certificates). Subscribers are required to create key pairs securely, maintain control of their private keys and notify the Certifying Authority immediately in case of compromise. According to this chapter, if a subscriber accepts a Digital Signature Certificate, it certifies that the subscriber holds the private key corresponding to the public key in the certificate, and that all the information in the certificate that is withing the knowledge of the subscriber is true.

\section*{Chapter IX - Penalties, Compensation and Adjudication}

Chapter IX deals with penalties and compensation for cybercrimes like DOS attacks, introduction of viruses, data breaches, etc. A company is liable to paying compensation for the failure to protect user data. This chapter also enables the Central Government to appoint adjudicating officers to hear such cases and award compensation or penalties based on factors like financial gain from the offense or repetitive nature.

\section*{Chapter X - The Appellate Tribunal}

This chapter of the IT Act, appoints the Telecom Disputes Settlement and Appellate Tribunal as the tribunal to hear appeals against orders passed by the adjudicating officers or the controller. The chapter also defines the tribunal's jurisdiction, procedures, powers, and appeal timelines.

\section*{Chapter XI - Offences}

\par Chapter XI lists various cybercrimes and along with their criminal penalties such as imprisonment and fines. Cybercrimes include tampering with computer source code, hacking, identity theft, sending offensive messages, cheating by impersonation, violation of privacy, cyber terrorism, publishing obscene content, including obscene content involving children, etc.

\section*{Chapter XII - Intermediaries not to be liable in certain cases}

\par Chapter XII of the IT Act states the legal exemption of intermediaries such as internet service providers, social media platforms, search engines, etc. where users can create content. It shields them from liability if they merely act as conduits for information and follow government-prescribed due diligence, including prompt removal of unlawful content when notified.

\section*{Chapter XIIA - Examiner of Electronic Evidence}

\par According to Chapter XIIA, the Central Government has the power to appoint Examiners of Electronic Evidence to assist courts and authorities in verifying the authenticity and content of digital records. These examiners provide technical validation of electronic evidence during investigations and legal proceedings.

\section*{Chapter XIII - Miscellaneous}

\par The final chapter of the IT Act addresses various legal and administrative matters. It empowers police officers and other authorized personnel to search, seize, and arrest in cases involving cyber offenses. It designates government officials under the Act as public servants and offers them legal protection for actions taken in good faith. The chapter establishes the Act's overriding authority over conflicting laws and extends it applicability to modern financial tools like electronic cheques and truncated cheques.


\end{document}
